\documentclass[11pt,a4paper]{article}
\usepackage[margin=1in]{geometry}

\usepackage[scaled=0.95]{helvet}
\renewcommand{\familydefault}{\sfdefault}

\usepackage{xcolor}
\usepackage{enumitem}
\usepackage{amsmath}
\usepackage{titlesec}
\usepackage[colorlinks=true,allcolors=blue]{hyperref}

\titleformat{\section}{\large\bfseries}{\thesection.}{0.6em}{}
\titlespacing*{\section}{0pt}{16pt}{6pt}

\newlist{qa}{enumerate}{1}
\setlist[qa]{label=\textbf{\arabic*.}, leftmargin=2.2em, itemsep=8pt, topsep=6pt}

\newcommand{\C}[1]{\textbf{#1}\par\vspace{2pt}}
\newcommand{\R}[1]{\textbf{Response.}~#1}

\begin{document}

\begin{center}
{\Large\bfseries Response to Reviewers}\\[4pt]
% {\normalsize Manuscript: \emph{Learning Task-Relevant Latent Residuals for Robot Manipulation}}\\[2pt]
% {\normalsize (previously submitted to IEEE/RSJ IROS 2026 (Submission \#2433).}
\end{center}

\vspace{6pt}
\noindent
Please find enclosed the revised version of the manuscript ``Learning Task-Relevant Latent Residuals for Robot Manipulation'', previously submitted to IEEE/RSJ IROS 2026 (Submission \#2433). We thank the Associate Editor and the reviewers for their constructive comments. We have substantially revised our manuscript based on the reviewers' comments, and the main changes are:

\begin{itemize}[leftmargin=1.4em,itemsep=2pt,topsep=4pt]
  \item \textbf{The asymmetric prediction mechanism is re-described.} The two streams are now
        presented as \emph{one} shared predictor $P$ with shared future queries $Q$, evaluated
        under two memory-access conditions (Eqs.~1--2), and the section closes with an explicit
        statement of where gradients stop and which modules are frozen versus trainable
        (Sec.~III-A).
  \item \textbf{The stream retained at deployment is corrected.} It is the \emph{intent} stream,
        not the full-information stream, that is kept after pre-training. Sec.~III-B now specifies
        the output structure of the retained predictor, the read-outs used by the two policy
        classes, and the closed-loop execution frequency --- including what counts as one control
        step in each of the two action spaces.
  \item \textbf{A second simulation benchmark.} Meta-World MT10 is added (Sec.~IV-C,
        Table~III): ten heterogeneous tasks covered by a single language-conditioned policy
        trained by offline behavioural cloning. Both simulation benchmarks are now averaged over
        three random seeds, and the real-robot evaluation has been extended from 5 to 15 trials
        per subtask.
  \item \textbf{Experimental protocols are specified.} The revised experiments provide the
        applicable camera views, image resolution, proprioceptive inputs, frame stacking, depth
        usage, object randomisation and separation between demonstration and evaluation seeds.
  \item \textbf{Relation to latent action models.} Sec.~II-B now contrasts TRAIL with latent action
        pre-training.
  \item \textbf{Figures and wording.} The pre-training architecture figure has been redrawn with a
        frozen/trainable legend, shared-weight markers, an explicit stop-gradient marker and the
        two memory-access conditions drawn separately; the affordance visualisation has been
        replaced; and the introduction's claim about discarded predictive decoders no longer
        cites R3M as an instance of it.
  \item \textbf{Supplementary videos.} The project webpage linked from the manuscript now labels
        the playback speed of every video group, distinguishing accelerated, real-time and
        slowed-down footage.
\end{itemize}

\vspace{2pt}
\noindent
\textbf{A note on naming.} The method was called \emph{LARM} in the IROS submission and is called
\emph{TRAIL} (\textbf{T}ask-\textbf{R}elevant residuals as an \textbf{A}ction prior for robot
man\textbf{I}pu\textbf{L}ation) in this version; the title was shortened accordingly. The method
itself is unchanged apart from the corrections described above. We use the new name throughout this
letter.

\vspace{2pt}
\noindent
This Statement of Changes includes our point-by-point response to each comment of the Associate
Editor and the reviewers. Section, equation, figure and table numbers refer to the revised
manuscript unless stated otherwise.

% ------------------------------------------------------------------
\section{Associate Editor}

\begin{qa}

\item \C{The idea is novel and performs reasonably well, but the submission lacks clarity.
In particular, the paper states that $F_{\mathrm{intent}}$ is discarded after training, while
none of the losses affecting $F_{\mathrm{full}}$ appears to depend on $F_{\mathrm{intent}}$.
Should the stop-gradient in Eq.~(5) instead be applied to $F_{\mathrm{intent}}$? There may also
be coupling through the text backbone, making the method difficult to assess.}
\R{We thank the Associate Editor for identifying this ambiguity. The previous manuscript
incorrectly stated which stream is retained and did not make the parameter sharing or gradient
flow sufficiently explicit. In the revision, Sec.~III-A and Fig.~2 describe a \emph{single}
predictor $P$ with shared future queries $Q$, evaluated twice under different memory-access
conditions. The full-information pass receives
$[V'_{\texttt{cls}},V'_{\texttt{patch}},L']$, whereas the intent pass receives
$[V'_{\texttt{cls}},L']$; they are two passes through shared parameters, not two independently
trained predictors.

The stop-gradient in the consistency loss is intentionally applied to the full-information
prediction:
\[
\mathcal{L}_{\mathrm{cons}}=
\operatorname{MSE}\!\left(\hat z^{\mathrm{int}}_{t+k},
\operatorname{sg}\!\left(\hat z^{\mathrm{full}}_{t+k}\right)\right).
\]
This term treats the patch-informed full prediction as a fixed target and updates the intent pass.
It does not stop $\mathcal{L}_{\mathrm{full}}$ from training the shared predictor through the
full-information pass; $\mathcal{L}_{\mathrm{int}}$ likewise trains it through the intent pass.
Applying stop-gradient to the intent prediction instead would reverse the intended direction of
the consistency transfer. There is no gradient coupling through the text backbone because
DistilBERT is frozen; the current-frame visual encoder, language projection, fusion module,
future queries and shared predictor are trainable. These details are now stated after Eqs.~(3)--(6)
and visualised in Fig.~2.

Finally, the corrected deployment description in Secs.~III-A--III-B states that the
full-information access condition is used only during pre-training. The \emph{intent} stream is
retained, predicts the future token sequence, and supplies the token-wise residual $\Delta_z$ to
the downstream policy.}

\item \C{The video is a good teaser for the paper.}
\R{We thank the Associate Editor for the positive feedback. We have expanded the experimental
presentation on the project webpage and added explicit playback-speed labels to all video groups.}

\end{qa}

% ------------------------------------------------------------------
\section{Reviewer 8}

\begin{qa}

\item \C{This paper should clarify the similarities, differences, and advantages compared to
latent action training. Prior works [1, 2] have utilised a latent action representation for policy
training with action-free video datasets. However, this work does not train this latent
information; how could this method make the model learn the dynamics information?}
\R{Thanks for the suggestion. We have added an explicit comparison in Sec.~II-B, marked in red.
\emph{Similarity:} both approaches exploit temporal transitions in action-free video.
\emph{Difference:} latent action models infer a low-capacity latent action from the current and
future frames, use it to reconstruct future-frame pixels, and subsequently map that latent action
to executable robot actions. TRAIL does not learn an action variable. It predicts the frozen visual
representation of a future frame from the current frame and instruction, then supplies the
observation-space residual between the current and predicted future representations to the policy.
\emph{Advantage:} TRAIL operates in latent feature space rather than reconstructing pixels and does
not require a separate latent-to-executable-action mapping stage. The answer to how dynamics are
learned without a latent action is therefore the future-latent prediction objective itself: to
minimise this objective, the predictor must encode the task-relevant transition implied by the
current scene and language instruction. The downstream imitation-learning policy then learns the
robot action from labelled demonstrations while using that transition as context.}

\item \C{The Franka Kitchen simulation is too simple to validate the effectiveness of this method.
The authors should provide more simulation results.}
\R{Thanks for your suggestion. We have conducted another experiment on a harder simulation benchmark that consists of 10 different robotic manipulation tasks. Sec.~IV-C
and Table~III report the results on Meta-World MT10: ten tasks with distinct scenes, objects and motion primitives, covered by a \emph{single} language-conditioned policy trained by imitation learning. All backbones are frozen and followed by the same pooled action decoder, so the representation supplied to the decoder is the only quantity that varies across rows. We report mean success over three random seeds with 50 closed-loop episodes per task and seed, using demonstration splits, seeds, camera and proprioceptive inputs and evaluation configurations that are shared across
methods. Franka Kitchen results are now also reported over three random seeds and two viewpoints.

TRAIL reaches 78.2\% against 70.0\% for the strongest baseline. The performance gap mainly comes from the two tasks including moving objects, Pick-Place (40.0\% versus at most 18.0\% for any baseline) and Push (70.0\% versus at most 38.0\%), which require the policy that considers the object movement. We also show Fig.~5 that decomposes the failures of these two tasks by the furthest achieved stage. The results show that TRAIL reduces failures before object contact and shifts the remaining Pick-Place failures towards the later placement stage.}

\item \C{The hyper-parameters are strange. It does not make sense to select $k=96$ to provide the target future state, since it is a long time for precise robotic manipulation tasks.}
\R{We agree that the previous description could be misread as a robot control horizon. Here,
$k=96$ is only the frame gap used to construct pre-training pairs from Kinetics-400; it does not
mean that a robot executes 96 steps open loop. At the dataset's $25$--$30$\,fps, this gap spans
roughly $3$--$4$ seconds of human activity and is chosen to capture a meaningful state transition
rather than nearly identical adjacent frames. At deployment, the retained intent stream is
re-evaluated from a fresh observation at every policy decision. The residual is used as a
directional feature, while the action decoder learns the action magnitude and time scale from
labelled robot demonstrations. We now state this distinction explicitly in Sec.~IV-A, where $k$
is introduced, and Sec.~III-B, where closed-loop deployment is described.}

\end{qa}

% ------------------------------------------------------------------
\section{Reviewer 9}

\begin{qa}

\item \C{The notation around the \texttt{[CLS]} token is ambiguous. Section III-A refers to a
\texttt{[CLS]} token without clearly specifying from which modality it originates. Eq.~(4) uses
$V_{\mathrm{cls}}'$, which suggests the visual-side CLS token, but this should be stated
explicitly, especially since the language encoder may also contain special tokens.}
\R{Thanks for your comment. We have clarified the notation. Sec.~III-A now introduces the visual sequence as one class token $V_{\texttt{cls}}$ and $196$ patch tokens $V_{\texttt{patch}}$, and the instruction separately as the language-token sequence $L$. $V'_{\texttt{cls}}$, $V'_{\texttt{patch}}$ and $L'$ denote the same quantities after bidirectional cross-attention. We have also updated Eqs.~(1)--(2), Fig.~2 and the corresponding paragraphs in Sec.~III-A.}

\item \C{The paper states that the intent predictor is discarded at deployment, while the
full-information predictor is retained to produce $z_{t+1}$. However, it remains unclear what
output structure this retained predictor produces during control. In particular, since Sec.~III-C
uses a ``\texttt{[CLS]} residual'' for 1-D control, could the authors clarify how the
\texttt{[CLS]} token is obtained at deployment?}
\R{We apologise for the confusion. The previous sentence was incorrect: the \emph{intent} stream
is retained after pre-training, and the full-information access condition is removed. The retained
predictor emits a complete future token sequence aligned with the visual target, so Eq.~(7) forms
$\Delta_z$ token by token. For continuous-control benchmarks, $z_t$ and $\Delta_z$ are concatenated
with proprioception and passed to the action decoder. For the real-robot tabletop tasks, their patch
tokens are reshaped into spatial feature maps and fused to predict language-conditioned affordance
maps. Thus, the spatial output is predicted directly from the future queries rather than recovered
from a single \texttt{[CLS]} token. We have corrected and clarified this in Secs.~III-A--III-B and
the Fig.~2 caption.}

\item \C{Currently, Figure 2 is difficult to understand. The text backbone is not marked as either
trainable or frozen, the figure lacks arrows indicating how information flows through the
bidirectional attention modules, and the caption terminology does not appear in the figure.}
\R{Figure~2 has been redrawn, including an explicit legend with four entries: frozen,
shared weights, stop-gradient, and bidirectional cross-attention, and every module is marked
accordingly: the text backbone (DistilBERT) and the DINOv3 teacher are frozen, and all
remaining modules are trainable during pre-training.}

\item \C{According to Eq.~(4), $F_{\mathrm{intent}}$ uses $V'_{\mathrm{cls}}$ and $L'$, whereas the
figure seems to suggest that it takes only the CLS token as input.}
\R{We have aligned Eqs.~(1)--(2) and Fig.~2. They now show one shared predictor $P$ and shared
future queries $Q$ under two memory-access conditions:
$\hat z^{\mathrm{full}}_{t+k}=P(Q\mid[V'_{\texttt{cls}},V'_{\texttt{patch}},L'])$ and
$\hat z^{\mathrm{int}}_{t+k}=P(Q\mid[V'_{\texttt{cls}},L'])$. The revised figure explicitly draws
the class and language tokens entering the intent pass, and the class, patch and language tokens
entering the full-information pass.}

\item \C{In both simulation and real-world experiments, the visual inputs are not clearly
specified. Based on Fig.~6 the real-robot setup appears to use a RealSense RGB-D camera mounted
near the end effector, but the observations shown as ``Obs'' in Fig.~3 do not seem fully consistent
with that viewpoint. Please clarify the image inputs and camera viewpoints used in both settings.}
\R{We have added the details for each experiment setting.

\emph{Meta-World} (Sec.~IV-C, paragraph ``Visual input''): the \texttt{corner2} third-person camera
rendered at $224\times224$; the policy receives a single RGB frame together with a $4$-dimensional
proprioceptive state (end-effector position and gripper aperture); no wrist camera, no depth
channel and no frame stacking.

\emph{Franka Kitchen} (Sec.~IV-B): the benchmark's two standard viewpoints, with results averaged
over both, under the same frozen-backbone behavioural-cloning protocol as all baselines.

\emph{Real robot} (Sec.~IV-D and the Fig.~7 caption): a wrist-mounted Intel RealSense observing a
$60\,\mathrm{cm}\times30\,\mathrm{cm}$ workspace. The policy uses RGB only; the depth stream is not
used, and we have corrected the text, which previously described the sensor as ``RGB-D'' in a way
that implied depth was an input.

On the inconsistency the reviewer commented: The two views are in fact the same sensor.
The wrist-mounted camera is not read from arbitrary poses along the trajectory. The arm returns to a fixed observation pose before each decision, and the image the policy consumes (the one shown as ``Obs'') is always captured from that canonical pose. What looks like a third-person overhead view is therefore the wrist camera at the observation pose, not a second, externally mounted camera. We have added this description to clarify the observation protocol in Secs.~III-B and IV-D.}

\item \C{The paper does not explain how occlusion is handled in the real-robot setup. If only a
wrist-mounted RGB-D camera is used, it is unclear how the policy can reliably observe the placement
region after grasping, when the gripper or grasped object may block the workspace.}

\R{For the real-robot tabletop tasks, the policy does not observe the workspace while an
object is held. The action space is a complete \emph{pick-and-place primitive}, rather than a
low-level command: from one initial observation, the tabletop read-out reshapes the patch tokens of
$z_t$ and $\Delta_z$ into spatial maps and predicts both picking and placing poses through
language-conditioned affordance maps. The primitive then executes without re-observation. Thus,
possible occlusion after grasping cannot affect the placement decision, which has already been made
from the initial unobstructed observation. We have made this action granularity explicit in
Sec.~III-B.}


\item \C{The deployment-time role of the predictor is also unclear. It is unclear whether the
predictor is run once at the beginning of an episode, periodically, or at every control step in a
closed-loop manner.}
\R{The predictor is executed once per \emph{policy decision}, always from a fresh observation, in
closed loop. It is never rolled out over multiple future steps and never executed only once per
episode. What the submitted version failed to make explicit is that a ``decision'' has a different
granularity in the two action spaces TRAIL is evaluated in, because the two use different
read-outs. We have added this to the deployment paragraph of Sec.~III-B, which now states that a
control step denotes one policy decision whose granularity follows the action space:

\begin{itemize}[leftmargin=1.4em,itemsep=3pt,topsep=4pt]
  \item \emph{Simulation (Franka Kitchen, Meta-World).} The policy emits low-level continuous
        commands, so one decision is one control step. The cycle --- observe, encode $z_t$, run the
        retained intent stream once to get $\hat z^{\mathrm{int}}_{t+k}$, form $\Delta_z$, predict
        an action, move --- runs at the environment's control frequency, and the predictor is
        re-evaluated on every new observation.
  \item \emph{Real robot.} The policy emits a pick-and-place primitive, so one decision is one
        primitive. The arm sits at its fixed observation pose, one observation is encoded, the
        intent stream runs once, and the resulting affordance prediction supplies both endpoints of
        the motion; the primitive then executes without re-observation, after which the arm returns
        to the observation pose and the next decision is made from a new observation. This is
        precisely why the occlusion concern in comment~6 does not arise: no observation is ever
        taken while an object is held.
\end{itemize}

In both cases the loop is closed on a fresh observation at every decision, the predictor is
evaluated exactly once per decision, and the pre-training offset $k$ plays no role --- it does not
set the robot's action frequency, which the downstream decoder learns from the labelled
demonstrations. Fig.~1 (bottom) draws this loop.}

\item \C{The experimental protocol lacks important details on scene initialisation, and only five
evaluation trials per task seem insufficient to convincingly demonstrate the effectiveness and
robustness of the method.}
\R{We accept both points and have addressed them with additional experiments rather than with
qualifications.

\emph{Trial counts.} In simulation, every reported number is now a mean over three random seeds
with 50 closed-loop episodes per task and seed --- 150 episodes per task on MT10, and the Franka
Kitchen numbers are likewise averaged over three seeds and two viewpoints. On the real robot we
have tripled the evaluation from 5 to 15 independent trials per subtask, i.e.\ 105 trials per
method across the seven subtasks. The conclusions are unchanged: TRAIL reaches 90\% overall against
roughly 80\% for the strongest baseline, LaVA-Man (Fig.~6).

\emph{Scene initialisation in simulation.} Object positions are randomised by the benchmark at
every reset, and the evaluation seeds control that randomisation. Crucially, the evaluation seeds
are disjoint from all demonstration-collection seeds, so no policy is scored on an object
configuration it saw during data collection. All methods are given the same demonstration splits,
the same seeds, the same camera and proprioceptive inputs, the same action decoder and the same
evaluation configurations, so the comparison is paired.

\emph{Scene initialisation on the real robot.} At the start of each trial the objects are placed by
hand at randomised positions and orientations within the $60\,\mathrm{cm}\times30\,\mathrm{cm}$
workspace, so no trial repeats a memorised layout. All four methods --- CLIPort, MPI, LaVA-Man and
TRAIL --- are evaluated on the same set of initial configurations, so the 15 trials per subtask
form a paired comparison rather than four independently sampled ones. Together with the strict
exclusion of the five held-out blocks and five held-out objects from the training data, this is
what makes the margins on ``Packing unseen blocks'' and ``Packing unseen objects'' (Fig.~6)
attributable to the representation rather than to a more favourable draw of initial states.}

\item \C{Minor: it may not be appropriate to include R3M as reference [2] in support of the claim
that existing models discard the predictive decoder. R3M's dynamics learning is driven primarily by
time-contrastive learning, so discarding its auxiliary head is not equivalent to discarding the
dynamics information that predictive learning is intended to capture. The wording should be revised
or the citation removed.}
\R{We agree with the reviewer's distinction and have rewritten the sentence to make it. The
introduction now separates the two cases explicitly: predictive decoders \emph{may be removed
before downstream learning} in masked- and future-prediction pipelines (cited to MVP and masked
visual pre-training), \emph{while} objectives such as time-contrastive learning (cited to R3M)
leave temporal information only implicitly inside a static embedding. R3M is therefore no longer
offered as an instance of discarding a decoder, but as an instance of the second, distinct
situation --- which, as the reviewer notes, still leaves the downstream policy to recover that
information. The sentence that follows states the resulting challenge in those terms: not how to
learn future prediction, but how to retain the task-relevant part of it as an explicit policy
input.}

\item \C{The demonstration in the video is too fast. Also, are all of the demonstrations shown in
real time? If not, it would be helpful to indicate the playback speed used for each one.}
\R{We agree and have added explicit playback-speed labels to the supplementary project webpage
linked in the manuscript footnote. The physical UR5 videos are shown at $1.75\times$ speed, the
Franka Kitchen videos at $1\times$ (real time), and the Meta-World videos at $0.5\times$ speed. The
page displays these values in a global legend and beside the corresponding video groups, so it is
clear which footage is accelerated, real-time or slowed down.}

\end{qa}

% ------------------------------------------------------------------
\section{Reviewer 13}

\begin{qa}

\item \C{The reviewer's assessment is positive on both the conceptual contribution --- retaining
the predictor and computing explicit latent residuals rather than discarding the predictive decoder
--- and the experimental validation, and requests no specific revisions.}
\R{We thank the reviewer for their time in reviewing our paper. In this resubmission, we
have kept our methodology while improving presentation of the paper as well as adding additional simulation experiments and analysis. }

\end{qa}

\end{document}
